% !TeX root = ../papers/nc-p3.tex
\section{量子境界に退化する全窓の十八座標・四極大軌跡}
\label{sec:quantum-attraction-fiber}
\begin{theorem}\label{thm:quantum-attraction-fiber}
前節の階数12の量子類 $w$ を軌道閉包に含む窓の添字保持同型類は、次の十八座標で尽くされる。
四成分の和が零の指標格子で
\[
 \mathcal W=S_4(2,0,-1,-1)\ \cup\ S_4(1,1,-1,-1)
\]
とする。$z=(z_u)_{u\in\mathcal W}$ の非零台がある開半空間
$\langle u,\lambda\rangle>0$ に入るもの（空台を含む）を取り、
$T=(k^\times)^4/k^\times$ の $z_u\mapsto t^u z_u$ と四添字の偶置換で割る。
この対応は横断切片と固有線の基底を選んだ上での同型類の全単射である。
全点が標準Koszul左右AS正則である。極大な台は偶置換まで四つあり、
同型を除いた一般次元は $8,6,6,6$、これらの全台上で安定化群は自明、最小周期は4である。
\end{theorem}
\begin{proof}[横断する十八文字]
$q_{0j}=1$, $(q_{12},q_{13},q_{23})=(i,-i,i)$ と取る。
行積は全て1で、射影テンソルの巡回は量子体積を固定する。
基点は射影的にpolystableである。複素数上の正規直交基底では24個の量子交代係数が
全て絶対値1で、各一位置の四切片の台は互いに交わらずノルムが等しい。
従って定理\ref{thm:walsh-presentation-space}と同じ最小ノルム判定が使える。
四三角積による同型判定から、四座標線の置換で基点を保つものは $A_4$ である
（前節の軌道長2）。位置0での行列が与えられると、他位置の行列は六関係から射影的に一意である。
各許された置換は位置0の純置換行列に取れ、必要な位置別倍率は三角積判定の対角twistで得る。
純置換の積も位置0で決まるので、全射影安定化群は $H=T\rtimes A_4$ である。

文字符号をテンソル語の四文字の個数から $(1,1,1,1)$ を引いたものとする。
\path{computations/verify_quantum_attraction_fiber.py} は前節の接方程式をこの次数で分ける。
文字の型が $(3,-1,-1,-1),(2,0,-1,-1),(1,1,-1,-1),(1,0,0,-1),0$ の順に、
一文字あたりのアフィン接次元の上界は $0,1,1,4,13$ である。
合計79は前節の滑らかな76次元と三量子方向の和に等しく、個々に等号である。
$Z_5$ では零文字だけが3減る。四基底軌道と全体倍率の接では
$(1,0,0,-1)$ の各文字に四位置から四方向があり、零文字に十方向がある。
実際、射影安定化Lie環は対角三次元である。従って横断表現が定理の十八線となる。
十二本の軌道で $A_4$ は単純推移的であり、基底を置換して働く形にできる。
六本の軌道では $(1,1,-1,-1)$ の安定化元 $(01)(23)$ を調べる。
位置 $n$ で文字 $j$ を $i^{n d_j}$ 倍して置換し、$d=(0,0,3,1)$ と取ると
体積とこの変形線の双方の倍率が $-1$ である。従って射影接への倍率は1であり、
この六本にも純置換基底が取れる。計算は法97で行える。固有値は符号または四乗根で、
単根が剰余でも別々なので、得る指標は標数零で同じである。
\end{proof}

\begin{proof}[零ファイバーの全軌道]
閉軌道のLunaのエタール切片定理を、指数の巡回移動も加えた有限拡大の群で適用する。
射影空間では基点で非零の不変斉次式のアフィン開集合を使う。
窓の空間の切片はここで滑らかな $Z_5$ に取り、その切片から上の横断接表現へも
安定化群同変なエタール写像を取る。例えば、安定化群不変な局所アフィン座標で
軌道と横断方向の不変補空間を取り、横断接への射影を用いる。
いずれも基点の周囲でstrongly étale、すなわち不変商の射と元の射との正方形を
ファイバー積にできる（Lunaのエタール切片定理と基本補題）。
二つの商の基点ファイバーへの制限は従って同型である。
特にこれは完備局所環だけでなく、線型表現の零ファイバーとテンソル側の基点ファイバーの
全幾何軌道の対応を与える。飽和開集合は基点ファイバーの全体を含む。
巡回を加えても基点の元の四基底軌道は同じで、安定化群の元の部分群に戻して
添字保持同型だけを取れる。なお周囲の階数スキームでも $w$ を軌道閉包に含む点は、
閉軌道への一径数判定と前節の流入の定理により $Z_5$ に入る。

有限拡大 $T\rtimes A_4$ の線型零ファイバーは $T$ のものと同じである。
トーラスの対角作用では、原点へ退化する条件は非零台が開半空間に入ることである。
実際、その条件を満たす有理余指標を整数倍すれば直接退化でき、
満たさなければ台上で非零の正次数不変単項式がある。
その有限軌道の値を区別する不変量も取れる。成分表現を射影商でさらに商にした一点を
同型類空間とするのではなく、この零ファイバー内の個々の軌道を数えている。
\end{proof}

\begin{proof}[全特殊値の同型、周期と退化]
$\mathcal W$ を辞書式に並べて $0,\ldots,17$ と番号を付ける。
共通の台 $S$ の二点のトーラス同型は、$\mathbb Z^S\to X^*(T)$ の整数核の各基底元 $m$ で
$\prod_{j\in S} z_j^{m_j}$ が一致することと同値である。これは代数閉体上のトーラスの像の方程式であり、
飽和していない指標像でも成立する。先に偶置換全てを試せば台の異なる場合も含む。
一位置巡回は上記順の各座標に $i^{b_j}$ を掛け、向きを一方に取ると
\[
 (b_j)=(3,2,1,1,3,2,2,3,1,0,0,0,0,0,0,0,0,0).
\]
これは各一元接核で語を巡回して体積の符号を除く前掲コードの計算による。
従って $r=1,2,4$ の順に $(i^{rb_j}z_j)$ と $z$ の上の同型検査を行えば全台と全特殊値の最小周期が決まる。
退化も、偶置換後にある余指標が残す台 $S'$ で零、$S\setminus S'$ で正であり、
残す座標が上のトーラス同型条件を満たすことと同値である。
これは対角トーラス軌道閉包の面による分解で、各面上の非零座標は一軌道となる。

極大台の列挙では第四余指標座標を零に取って三空間の中心超平面配置を使う。
各室の符号は、二枚の独立な壁の交線の半直線、その壁内の隣接方向、壁外の方向の順に
無限小変位を取って尽くせる。実際、各室の閉錐は尖っており、その稜と稜を含む面を一組選べる。
壁上の台は周囲の室の台に含まれる。全168室の正台のうち包含極大なのは十一座標が12個、
九座標が36個で、偶置換を除いた代表は
\[
\begin{gathered}
 \{0,1,2,3,4,5,6,7,8,9,10\},\qquad
 \{0,1,2,3,4,5,7,9,12\},\\
 \{0,1,2,3,4,5,9,10,12\},\qquad
 \{0,1,2,3,4,6,8,10,13\}.
\end{gathered}
\]
各代表は十二個の偶置換全てで別の台となり、指標行列の三Smith因子は全て1である。
従って台全体が非零なら安定化群は自明である。各整数核への $b$ の値と4の最大公約数は1で、
台を保つ置換もないので全台上の周期は4である。台の大きさから3を引くと表示した次元を得る。
\end{proof}
